% Capítulo 1: Introdução
\chapter{Introdução}\label{chp:introducao}

O sensoriamento remoto hiperespectral estabeleceu-se como tecnologia de imageamento 
avançada que revoluciona métodos de caracterização de materiais, superfícies e 
fenômenos naturais através da aquisição simultânea de informações espaciais e 
espectrais de alta resolução. Os sistemas hiperespectrais capturam informações 
detalhadas sobre composição química e física de objetos através da aquisição 
de centenas de bandas espectrais estreitas (tipicamente 100-300 bandas), cada 
uma com largura espectral de 2-10 nm, cobrindo faixas específicas do espectro 
eletromagnético na região VNIR-SWIR (400-2500 nm) \cite{Lorenzzetti2015}. 

Esta capacidade contrasta significativamente com sistemas multiespectrais 
convencionais, como Landsat TM5 (6-12 bandas com largura de 60-200 nm), que 
operam com bandas espectrais relativamente largas e espaçadas. Os sensores 
hiperespectrais oferecem resolução espectral superior (FWHM $\leq$ 10 nm), 
permitindo identificação de feições de absorção estreitas e assinaturas 
espectrais diagnósticas imperceptíveis em sistemas tradicionais. A razão 
sinal-ruído (SNR) típica de sistemas hiperespectrais modernos varia entre 
500:1 e 1000:1, proporcionando precisão suficiente para análises quantitativas 
de composição material \cite{Lorenzzetti2015}.

\begin{figure}[htbp]
    \centering
    \includegraphics[width=0.8\textwidth]{fig-cap01-01-principais-componentes-sensoreamento-remoto}
    \caption{Principais componentes de um sistema de sensoriamento remoto hiperespectral 
             embarcado: sensor de imageamento, unidade de processamento heterogênea 
             (FPGA/GPU), sistema de armazenamento e interface de comunicação para 
             aplicações em agricultura de precisão.}
    \label{fig:componentes-sistema-hiperespectral}
\end{figure}

Esta tecnologia habilita aplicações diversificadas em agricultura de precisão, 
proporcionando capacidades avançadas de detecção de estresse hídrico através 
da análise de feições de absorção da água (1450 nm e 1940 nm), identificação 
de deficiências nutricionais via índices de vegetação especializados (NDVI, 
NDRE, EVI, SAVI), e detecção precoce de infestações de pragas através de 
alterações nas assinaturas espectrais foliares. A \autoref{fig:componentes-sistema-hiperespectral} 
ilustra os componentes principais de um sistema hiperespectral embarcado integrado. \textcite{Adao2017} demonstram 
que UAVs emergiram como plataforma custo-efetiva para imageamento hiperespectral, 
substituindo satélites e aeronaves tripuladas em aplicações de monitoramento 
localizado, com classificadores Support Vector Machine (SVM) com kernel RBF 
alcançando Overall Accuracy (OA) de 92\% e coeficiente Kappa ($\kappa$) superior a 
0,90 utilizando apenas 10 amostras de treinamento por classe.

No monitoramento ambiental, os sistemas hiperespectrais permitem caracterização 
quantitativa de ecossistemas, detecção de poluentes através de suas assinaturas 
espectrais específicas, e acompanhamento temporal de mudanças climáticas com 
precisão radiométrica superior a 95\%. Em aplicações de defesa e segurança, 
a tecnologia detecta materiais camuflados explorando diferenças sutis nas 
propriedades espectrais, enquanto na prospecção mineral facilita identificação 
remota de depósitos através de feições de absorção diagnósticas de minerais 
específicos nas regiões VNIR e SWIR \cite{Baptista2019}.

O processamento de dados hiperespectrais --- que podem exceder centenas de 
megabytes por frame (tipicamente 200-500 MB para imagens 512×217×204 bandas) --- 
em plataformas embarcadas com restrições rigorosas de energia (10-20 W), peso 
($<$ 2 kg) e tamanho representa desafio computacional central da era moderna 
de sistemas autônomos. A complexidade computacional emerge da necessidade de 
processamento em tempo real (throughput $>$ 100 pixels/s) de algoritmos 
intensivos como Principal Component Analysis (PCA), classificação SVM com 
kernel RBF, e cálculo de índices de vegetação, mantendo simultaneamente 
precisão de classificação superior a 90\% (Overall Accuracy) e eficiência 
energética adequada para operação autônoma prolongada. Esta dissertação 
investiga sistematicamente este problema multidimensional, propondo abordagem 
metodológica rigorosa baseada em análise comparativa quantitativa e síntese 
de arquiteturas heterogêneas otimizadas.

Esta pesquisa aborda uma lacuna na literatura: ausência de frameworks para 
avaliação e otimização de técnicas de processamento hiperespectral em 
sistemas embarcados. Através de metodologia estruturada, investigamos 
trade-offs entre performance computacional, eficiência energética e precisão 
de classificação, buscando possíveis soluções arquiteturais para diferentes cenários 
de aplicação.

\section{Contexto e Motivação}\label{sec:contexto_motivacao}

O processamento de imagens hiperespectrais em sistemas embarcados representa 
desafio computacional multidimensional de complexidade exponencial. A 
complexidade fundamental emerge das características intrínsecas dos dados 
hiperespectrais, onde cada pixel constitui um vetor espectral de alta 
dimensionalidade ($\mathbb{R}^{200-300}$) contendo assinatura espectral 
composta por centenas de valores espectrais (tipicamente 150-300 bandas após 
filtragem atmosférica) que capturam propriedades biofísicas e bioquímicas 
específicas dos objetos observados através de processos de interação 
radiação-matéria \cite{Zhang2024}.

Esta riqueza informacional constitui simultaneamente a vantagem analítica 
e o desafio computacional central do imageamento hiperespectral. A dimensionalidade 
dos dados resulta em volumes típicos de 200-500 MB por frame para sistemas 
como Salinas Valley (512×217×204 bandas), exigindo processamento intensivo 
de operações matriciais de alta precisão ($\geq$ 32-bit floating point) para 
preservar características espectrais sutis. O desafio intensifica-se pela 
necessidade de processamento em tempo real (latência $<$ 100 ms por frame) 
sob restrições rigorosas de recursos computacionais, energéticos (TDP 10-20 W) 
e de memória (tipicamente $\leq$ 8 GB) característicos de plataformas embarcadas 
para UAVs.

O desafio computacional intensifica-se exponencialmente em aplicações de 
veículos aéreos não tripulados (UAVs) para imageamento hiperespectral, onde 
as condições operacionais introduzem complexidades adicionais de processamento. 
Imagens hiperespectrais capturadas por UAVs exibem variações espectrais 
pronunciadas (desvio padrão espectral 15-30\% superior) comparadas aos dados 
de plataformas espaciais estabilizadas, apresentando desafios computacionais 
específicos para algoritmos de classificação que assumem distribuições 
espectrais homogêneas \cite{Zhang2024}.

Esta complexidade surge das condições operacionais dinâmicas dos UAVs: 
variações de altitude (50-500 m), ângulos de observação variáveis (0-30° 
off-nadir), condições atmosféricas heterogêneas (vapor d'água, aerossóis), 
e variações de iluminação solar introduzem heterogeneidade espectral e 
radiométrica que requer correções adaptativas em tempo real. Sistemas 
tradicionais de processamento, projetados para dados espaciais estáveis, 
demonstram degradação de performance (redução 20-40\% na Overall Accuracy) 
quando aplicados diretamente a dados UAV sem pré-processamento especializado. 
A necessidade de correção geométrica dinâmica, calibração radiométrica 
adaptativa, e filtragem atmosférica em tempo real adiciona overhead 
computacional de 30-50\% ao pipeline de processamento.

O problema intensifica-se exponencialmente com requisitos rigorosos de 
processamento em tempo real para aplicações críticas. Sistemas hiperespectrais 
embarcados devem satisfazer simultaneamente múltiplos requisitos conflitantes: 
alta taxa de compressão (razões 10:1 a 50:1) devido aos volumes massivos de 
dados, processamento em tempo real com throughput mínimo de 100-330 pixels/s 
para evitar acúmulo de dados não processados (buffer overflow), latência 
determinística inferior a 100 ms para aplicações de tempo real, e complexidade 
computacional rigorosamente limitada (máximo 2.500 operações de ponto flutuante 
por pixel) devido às restrições de recursos das plataformas embarcadas 
\cite{Diaz2019}.

A análise quantitativa dos requisitos revela trade-offs fundamentais: 
algoritmos de redução dimensional como PCA requerem decomposição de matrizes 
de covariância ($d \times d$, onde $d \approx 200$ bandas), resultando em 
complexidade $O(d^3)$ que excede capacidades de processamento em tempo real 
para plataformas com TDP limitado. Classificadores SVM com kernel RBF, embora 
proporcionem alta precisão (OA $>$ 95\%), apresentam complexidade de inferência 
$O(n_{\text{sv}} \times d)$ onde $n_{\text{sv}}$ pode exceder milhares de 
support vectors, tornando-se computacionalmente proibitivos para processamento 
pixel-wise em tempo real.

A análise sistemática dos trabalhos existentes revela limitações fundamentais 
das técnicas tradicionais de processamento quando aplicadas ao contexto 
embarcado com restrições rigorosas. \textcite{Manolakis2002} demonstram que 
detectores estatísticos clássicos como Reed-Xiaoli (RX) e matched filters, 
embora teoricamente fundamentados em princípios de detecção ótima, assumem 
distribuições gaussianas multivariadas que frequentemente não se aplicam a 
cenários reais de dados hiperespectrais (teste Shapiro-Wilk com p-value $< 0{,}05$ 
para 60-80\% das bandas espectrais). A violação dessas premissas resulta em 
degradação significativa da performance (redução 15-25\% na taxa de detecção 
para false alarm rate fixo de $10^{-4}$).

\textcite{Nasrabadi2014} complementa esta análise identificando que detectores 
adaptativos como ACE (Adaptive Coherence Estimator) e GLRT (Generalized 
Likelihood Ratio Test) superam métodos de correlação simples em cenários 
não-gaussianos, mas requerem processamento computacional intensivo: ACE 
demanda inversão de matrizes de covariância $O(d^3)$ e GLRT necessita 
otimização iterativa com 50-100 iterações por pixel. Para dados Salinas 
Valley ($d = 204$ bandas), estas operações excedem 10$^6$ operações de 
ponto flutuante por pixel, inviabilizando processamento em tempo real em 
plataformas com TDP $\leq$ 20 W.

\textcite{Diaz2019} observam que, embora compressores com perdas alcancem 
performance satisfatória em termos de taxa de distorção (PSNR $>$ 40 dB 
para razões 20:1), eles são caracterizados por custos computacionais 
exponenciais (complexidade $O(n^2 \log n)$ para transformadas wavelets 2D), 
requisitos intensivos de memória (buffers $>$ 100 MB para processamento 
por blocos), e natureza não escalável com resolução espectral, impedindo 
seu uso em aplicações sob restrições rigorosas de latência ($< 100$ ms), 
energia ($< 2$ J/frame), e memória ($< 8$ GB), típicas da compressão 
embarcada em UAVs.

A demanda crescente por aplicações autônomas tem impulsionado o desenvolvimento 
acelerado de soluções baseadas em arquiteturas heterogêneas especializadas, 
que combinam diferentes tipos de processadores para otimização simultânea de 
throughput, latência e eficiência energética. O avanço dos UAVs, combinado 
com a miniaturização de equipamentos de imageamento hiperespectral (redução 
de peso de 50-80\% em sistemas como Specim FX10 e Resonon Pika L), tornou o 
sensoriamento remoto hiperespectral baseado em UAV uma tecnologia essencial 
para observação da Terra com resolução temporal e espacial sem precedentes 
\cite{Zhang2024}. 

No entanto, \textcite{Zhang2024} ressaltam que os sensores hiperespectrais 
baseados em UAV, embora ofereçam vantagens significativas em capturar 
informações espectrais de alta resolução de objetos terrestres (GSD de 
1-10 cm para altitudes 50-200 m), também apresentam desafios computacionais 
exponenciais relacionados ao processamento de dados de alta resolução 
espacial (até 4096×4096 pixels), temporal (até 60 fps), e espectral 
(200-400 bandas), resultando em taxas de dados de 1-10 GB/s que excedem 
capacidades de armazenamento e transmissão de plataformas embarcadas 
convencionais.

O cenário atual de pesquisa revela avanços significativos em técnicas 
individuais de otimização, com deep learning emergindo como paradigma 
dominante para classificação hiperespectral. Técnicas avançadas incluem 
redes neurais convolucionais (CNNs) com arquiteturas 3D que exploram 
correlações espectral-espaciais, métodos baseados em Transformer que 
aplicam mecanismos de atenção para seleção adaptativa de bandas, redes 
neurais recorrentes (RNNs) para modelagem sequencial de assinaturas 
espectrais, redes neurais de grafos convolucionais (GCNs) que capturam 
relações topológicas entre pixels, redes adversariais generativas (GANs) 
para aumento de dados e correção de domínio, e modelos híbridos que 
combinam múltiplas arquiteturas, demonstrando potencial superior na 
classificação de imagens hiperespectrais com Overall Accuracy típica 
de 85-98\% em datasets padronizados \cite{Zhang2024}.

\textcite{BioucasDias2012} demonstram que técnicas de unmixing espectral, 
incluindo abordagens geométricas (N-FINDR, VCA), estatísticas (ICA, NMF) 
e de regressão esparsa (SUnSAL, CLSUnSAL), são fundamentais para análise 
quantitativa de pixels mistos, com métodos spectral-spatial alcançando 
RMSE 15-30\% inferior comparado à análise apenas espectral. A integração 
de informações espaciais através de filtros morfológicos, regularização 
espacial, e análise de contexto local pode melhorar significativamente 
a precisão (ganhos 5-15\% na Overall Accuracy) e robustez da classificação, 
explorando características espaciais e espectrais simultaneamente através 
de arquiteturas end-to-end otimizadas.

Paralelamente, \textcite{Diaz2019} demonstram o potencial transformador das 
GPUs embarcadas para compressão hiperespectral em tempo real, implementando 
o compressor HyperLCA (Lossy Compression Algorithm for Hyperspectral Image 
System) otimizado para arquiteturas CUDA nos kits de desenvolvimento NVIDIA 
Jetson TK1 (192 cores CUDA, TDP 10 W) e TX2 (256 cores Pascal, TDP 7,5-15 W). 
Os resultados experimentais obtidos pelos autores demonstram que o compressor 
HyperLCA consegue processar consistentemente mais de 330 frames por segundo 
(throughput sustentado) para imagens 512×512×224 bandas, com razão de compressão 
20:1 e PSNR superior a 40 dB, verificando simultaneamente a eficácia do 
algoritmo HyperLCA e o alcance de performance em tempo real através das 
implementações CUDA otimizadas.

Este trabalho estabelece marco fundamental demonstrando como GPUs embarcadas 
de baixo consumo podem ser efetivamente utilizadas para acelerar algoritmos 
computacionalmente intensivos de processamento hiperespectral, aproveitando 
sua arquitetura massivamente paralela (centenas a milhares de cores) para 
processamento SIMD de operações matriciais típicas. A eficiência energética 
alcançada (>20 GOPS/W) supera implementações CPU tradicionais por fatores 
de 5-10x, demonstrando viabilidade de processamento hiperespectral embarcado 
em plataformas com severas restrições energéticas.

A importância crítica da eficiência energética em sistemas embarcados é 
demonstrada quantitativamente pelos resultados que mostram melhorias 
exponenciais através de arquiteturas heterogêneas otimizadas 
\cite{Vaithianathan2025}. Técnicas avançadas de co-design hardware-software, 
que integram otimizações algorítmicas com implementações arquiteturais 
especializadas, podem resultar em melhorias dramáticas de 21x na velocidade 
de throughput (de 15 para 330 fps) \cite{Hwang2011} e reduções extraordinárias 
de até 69,88x no consumo energético (de 140 W para 2 W) comparado a 
implementações GPU convencionais não-otimizadas \cite{ArucuIliev2025}.

Os resultados de co-design evidenciam o potencial transformador de abordagens 
integradas que exploram sinergias computacionais entre diferentes componentes 
arquiteturais: FPGAs proporcionam eficiência energética superior (>100 GOPS/W) 
para operações de ponto fixo e pipelines determinísticos, GPUs embarcadas 
oferecem throughput paralelo otimizado (>20 GOPS/W) para operações matriciais 
de ponto flutuante, e CPUs ARM fornecem flexibilidade de controle e 
gerenciamento de sistema com consumo moderado (5-10 GOPS/W). A combinação 
estratégica desses elementos em arquiteturas heterogêneas permite explorar 
o ponto ótimo de performance por watt para cada estágio específico do 
pipeline de processamento hiperespectral.

Contudo, apesar destes avanços pontuais, é identificado uma lacuna 
na literatura atual: ``Apesar destes avanços, as técnicas atuais de 
classificação de imagens hiperespectrais de UAV ainda enfrentam diversos 
desafios, incluindo processamento de dados de alta dimensionalidade, 
heterogeneidade espectral e variabilidade radiométrica entre diferentes 
cenas'' 
\cite{Zhang2024}. 
Os autores enfatizam que deep learning permanece largamente uma abordagem de 
``caixa preta'' para a maioria dos problemas, indicando a necessidade de 
abordagens mais sistemáticas e interpretáveis.

\section{Problema de Pesquisa}\label{sec:problema_pesquisa}

O problema central investigado nesta dissertação emerge da análise sistemática 
do estado da arte em processamento hiperespectral embarcado, revelando lacuna 
metodológica fundamental que permeia a literatura científica contemporânea e 
impede o desenvolvimento de soluções práticas otimizadas. Esta lacuna se 
manifesta através de três dimensões críticas identificadas na análise 
comparativa das técnicas de processamento hiperespectral desenvolvidas nas 
últimas décadas: fragmentação metodológica, ausência de frameworks de 
avaliação padronizados, e falta de síntese sistemática de trade-offs 
arquiteturais.

\textcite{Zhang2024} documentam que, ao longo das últimas décadas, as 
técnicas de classificação de imagens hiperespectrais de sensoriamento remoto 
por UAV alcançaram avanços significativos, oferecendo precisão analítica 
e detalhes espaciais sem precedentes para observação da Terra. O trabalho 
oferece revisão abrangente das aplicações e performances de métodos 
tradicionais de machine learning (SVM, Random Forest, k-NN) e deep learning 
(CNNs, RNNs, GCNs, Transformers) neste campo, resumindo as principais 
direções de pesquisa atuais e desafios técnicos específicos.

No entanto, os autores identificam fragmentação metodológica significativa 
na pesquisa: ``algoritmos tradicionais de machine learning podem alcançar 
certos níveis de precisão de classificação (Overall Accuracy 75-85\%), mas 
frequentemente têm dificuldades com a alta dimensionalidade (curse of 
dimensionality para $d > 100$ bandas) e natureza não-linear complexa dos 
dados hiperespectrais''. Esta limitação resulta em degradação sistemática 
de performance quando o número de bandas espectrais excede o número de 
amostras de treinamento (fenômeno de Hughes), problema particularmente 
crítico em aplicações embarcadas onde datasets de treinamento são limitados 
por restrições de armazenamento e processamento.

Esta fragmentação metodológica resulta em limitações críticas para o 
desenvolvimento de sistemas práticos otimizados. \textcite{Diaz2019} observam 
que a complexidade algorítmica dos métodos de compressão ($O(n^2 \log n)$ 
para wavelets 2D), combinada com volumes massivos de dados a serem comprimidos 
(200-500 MB/frame para imagens 512×217×204 bandas) e recursos computacionais 
rigorosamente limitados dos dispositivos de hardware embarcados (TDP $\leq$ 20 W, 
memória $\leq$ 8 GB, throughput $\geq$ 100 pixels/s), tornam a compressão 
em tempo real uma tarefa computacionalmente desafiadora que excede capacidades 
de plataformas convencionais.

Os autores enfatizam que esta complexidade não é meramente técnica, mas 
representa barreira fundamental à implementação prática de sistemas 
hiperespectrais em plataformas com restrições severas de recursos. A análise 
quantitativa revela que algoritmos estado-da-arte requerem 10$^4$-10$^6$ 
operações de ponto flutuante por pixel, resultando em demanda computacional 
total de 10$^{10}$-10$^{12}$ operações por frame, que excede capacidades 
de processamento em tempo real ($< 100$ ms/frame) para plataformas embarcadas 
típicas (1-10 GFLOPS sustentados).

A ausência crítica de metodologias sistemáticas para comparação objetiva 
entre diferentes abordagens arquiteturais cria ambiente de incerteza 
metodológica que impede significativamente o progresso científico da área. 
\textcite{Zhang2024} destacam que ``a performance de diferentes frameworks 
de redes neurais na classificação de imagens hiperespectrais de UAV varia 
significativamente, com métodos baseados em GCN e CNN geralmente superando 
abordagens baseadas em RNN (diferenças típicas de 5-15\% na Overall Accuracy)''.

No entanto, os autores observam que cada trabalho científico utiliza métricas 
de avaliação próprias (Overall Accuracy, Kappa, F1-score, RMSE aplicados 
inconsistentemente), datasets distintos (Indian Pines, Pavia University, 
Salinas Valley, Kennedy Space Center com diferentes pré-processamentos), 
condições experimentais particulares (razões de treinamento/teste variáveis 
de 10\% a 90\%, técnicas de validação cruzada inconsistentes), e plataformas 
de hardware heterogêneas (CPU Intel, GPU NVIDIA, FPGA Xilinx com diferentes 
gerações e configurações), impedindo comparações diretas quantitativas e 
síntese de conhecimento sistemática baseada em evidências.

O problema se agrava quando consideramos a predominância de validações em 
contextos controlados e simplificados. Vaithianathan observa que, embora 
avanços tenham sido alcançados em arquiteturas heterogêneas individuais, 
``questões relacionadas ao consumo de energia, escalabilidade, complexidade 
de programação, gerenciamento de recursos e segurança devem ser abordadas 
para criar sistemas robustos e eficientes''. O 
autor enfatiza que as complexidades da integração de múltiplos componentes, 
as interações entre subsistemas, e os desafios práticos de implementação 
permanecem inexplorados.

Díaz et al.\ contribuem para esta discussão documentando que os experimentos 
realizados em seu trabalho são ``orientados às necessidades impostas por uma 
aplicação específica de agricultura inteligente, embora todas as conclusões 
obtidas sejam extrapoláveis para outros campos nos quais imagens 
hiperespectrais sensoreadas remotamente devem ser comprimidas em tempo real''. 
Esta observação ressalta a necessidade de frameworks mais 
gerais que possam ser aplicados sistematicamente a diferentes domínios de 
aplicação.

Diante deste cenário multidimensional de lacunas metodológicas, o problema 
de pesquisa central pode ser formulado de forma precisa: Como desenvolver 
metodologia sistemática e reproduzível para investigação comparativa 
quantitativa de técnicas de otimização em processamento hiperespectral 
embarcado que permita identificação objetiva de arquiteturas ótimas através 
de análise rigorosa dos trade-offs entre performance computacional (throughput, 
latência), eficiência energética (GOPS/W, J/pixel), e precisão de classificação 
(Overall Accuracy, Kappa, F1-score) para diferentes cenários de aplicação 
em agricultura de precisão com requisitos específicos (tempo real, autonomia, 
precisão) e constraints operacionais distintos (TDP, memória, armazenamento)?

\section{Objetivos}\label{sec:objetivos}

Os objetivos desta dissertação foram estruturados metodologicamente para 
abordar sistematicamente as lacunas identificadas na literatura científica, 
respondendo diretamente às questões de pesquisa formuladas através de 
abordagem teórico-experimental rigorosa. O objetivo geral consiste em 
desenvolver metodologia sistemática e reproduzível para investigação 
comparativa quantitativa de técnicas de otimização em processamento 
hiperespectral embarcado, estabelecendo framework analítico integrado 
que permita quantificação objetiva e modelagem matemática de trade-offs 
multidimensionais entre performance computacional (throughput, latência), 
eficiência energética (GOPS/W, J/pixel), e precisão de classificação 
(Overall Accuracy, Kappa, F1-score) em arquiteturas heterogêneas FPGA-GPU 
aplicadas à agricultura de precisão.

Os objetivos específicos desdobram-se em cinco dimensões complementares 
com metas quantitativas precisas:

\textbf{Objetivo Específico 1 - Mapeamento Sistemático:} Realizar mapeamento 
abrangente de técnicas de otimização disponíveis na literatura científica 
(2019-2025), organizando-as segundo taxonomia estruturada que considere tipo 
de plataforma (FPGA Xilinx, GPU NVIDIA, CPU ARM), estágio de processamento 
(pré-processamento, redução dimensional, classificação, pós-processamento), 
métricas de performance específicas (throughput, latência, utilização de 
recursos, eficiência energética), e adequação quantitativa a cenários de 
agricultura de precisão (detecção de estresse hídrico, cálculo de índices 
de vegetação NDVI/NDRE/EVI/SAVI).

\textbf{Objetivo Específico 2 - Modelagem Matemática:} Desenvolver modelos 
matemáticos analíticos de performance para arquiteturas FPGA (Zynq-7020) e 
GPU (Xavier NX), estabelecendo framework teórico rigoroso que capture 
complexidade computacional ($O(n)$, $O(n^2)$, $O(n^3)$), utilização de 
recursos (DSP, BRAM, LUTs para FPGA; cores CUDA, memória, bandwidth para GPU), 
e padrões de consumo energético (potência instantânea, energia por pixel, 
eficiência GOPS/W) através de equações parametrizadas validadas experimentalmente.

\textbf{Objetivo Específico 3 - Ambiente de Simulação:} Implementar ambiente 
de simulação integrado baseado em ferramentas industriais (Vivado, 
GHDL, CUDA Toolkit) que permita avaliação comparativa sistemática 
sob condições experimentais controladas, integrando simuladores específicos 
de plataforma, dataset Salinas Valley (512×217×204 bandas) como referência 
padronizada, e protocolos de medição reprodutíveis com instrumentação de 
alta precisão (INA219 para FPGA, tegrastats para GPU).

\textbf{Objetivo Específico 4 - Análise Multi-critério:} Estabelecer 
metodologia de análise multi-critério baseada em teoria de decisão 
quantitativa, incorporando técnicas de otimização Pareto para identificação 
de soluções não-dominadas no espaço multidimensional de trade-offs, métodos 
de ponderação de critérios (AHP, TOPSIS) para diferentes perfis de aplicação, 
e análise de sensibilidade para validação da robustez das soluções identificadas.

\textbf{Objetivo Específico 5 - Síntese Arquitetural:} Propor arquitetura 
heterogênea otimizada baseada nos resultados da análise comparativa 
quantitativa, incluindo especificação arquitetural detalhada (particionamento 
de algoritmos, interfaces de comunicação, protocolos de sincronização), 
validação teórica através dos modelos matemáticos desenvolvidos, e estimativas 
de performance integrada (throughput $>$ 200 pixels/s, latência $<$ 50 ms, 
eficiência $>$ 30 GOPS/W, Overall Accuracy $>$ 90\%).

\section{Estrutura da Dissertação}\label{sec:estrutura}
Esta dissertação está organizada de modo a construir, de forma progressiva e
integrada, o conhecimento necessário para enfrentar o problema multidimensional
proposto. O Capítulo 2 apresenta o levantamento bibliográfico, no qual são
estabelecidos os fundamentos teóricos essenciais por meio de uma exposição
detalhada dos princípios físicos da espectroscopia de imageamento e da formação
de imagens espectrais, abordando aspectos como a interação radiação-matéria,
as regiões espectrais VNIR-SWIR (400-2500 nm), além de parâmetros relevantes
como SNR e FWHM. Nesse capítulo, realiza-se também uma análise sistemática das
tecnologias de câmeras hiperespectrais disponíveis, incluindo a discussão de
datasets amplamente utilizados, como Salinas Valley e Indian Pines, bem como
suas especificações técnicas. São ainda categorizadas, de maneira rigorosa,
as principais técnicas de processamento hiperespectral, como PCA, KPCA, EMCR,
SVM e métodos de classificação espectral, além de ser apresentada uma análise
comparativa das arquiteturas de processamento modernas, contemplando FPGAs
(Xilinx Zynq-7020), GPUs embarcadas (NVIDIA Jetson Xavier NX) e sistemas
heterogêneos, sempre com ênfase em métricas quantitativas de performance.

Na sequência, o Capítulo 3 é dedicado à metodologia, detalhando o delineamento
da pesquisa teórico-experimental e justificando cada escolha metodológica com
base em análise crítica da literatura. São descritos os protocolos de
implementação adotados para as plataformas FPGA (utilizando Vivado Design Suite e
GHDL) e GPU (com CUDA Toolkit), bem como os frameworks de modelagem
matemática empregados para análise da complexidade computacional. O capítulo
também aborda a construção dos ambientes de simulação integrados, que utilizam
o dataset Salinas Valley (512×217×204 bandas) como referência, e a aplicação
de instrumentação de precisão para medição energética, como INA219 e
tegrastats. O cronograma de execução estruturado em 10 meses (40 semanas) é
apresentado em seis fases principais, iniciando com aquisição de hardware e
culminando com publicação científica e defesa da dissertação. Por fim, são
formalizados matematicamente os procedimentos de análise multi-critério,
incluindo a utilização de métodos como otimização Pareto e as metodologias
AHP e TOPSIS para ponderação de objetivos, visando a identificação quantitativa
das soluções mais adequadas ao contexto estudado.